\begin{widefigure}[tbp]
    \includegraphics[width=\linewidth]{images/concepts.pdf}
    \caption{The structured coding scheme developed during Step 1 of the analysis to provide an overview of the discussions. Based on the interview structure, this scheme clusters recurring insights into seven topics that represent expert perspectives on multisensory biophilic design in indoor environments. Each topic contains subtopics and concepts discussed by experts.}
    \Description{Multi-Sensory Biophilic Design concept map diagram centred on "Multi-Sensory Biophilic Design. Branching clusters radiate outward, grouping terms around themes such as controlled versus idealized nature, technological enablers (e.g. smart glass, AI, generative models), participatory tools, values and societal dimensions (e.g. justice, inclusion, community), emotional and sensory effects (e.g. awe, surprise, restoration), and limitations (e.g. tokenism, ecological incoherence, economic constraints). The map includes interconnected terms and phrases reflecting design challenges, opportunities, and philosophical tensions in multisensory biophilic smart environments.}
    \label{fig:findings-findings-overview}
\end{widefigure}

\section{Findings: Topic-Level Clustering of Expert Responses}
\label{sec:ch4-findings-a}
This section presents the seven topics (TO1–TO7) derived from our topic-level clustering of expert interviews, offering a descriptive overview of how experts conceptualise biophilic design in smart buildings. An overview of the seven topics and their relationships is provided in \autoref{fig:findings-findings-overview}.

\paragraph{TO1 - A Deeper Connection to Nature:}
Experts spoke of designing for a ``deeper connection to nature'' in indoor spaces as a core philosophy underpinning biophilic design~\cite{kellert2008dimensions}. Beyond evident health benefits, they suggested that deeper connections could be fostered through existential moments of awe and humility, shifting perspectives from controlling nature to reflecting on humanity’s place within it. Experts stressed the importance of designing spaces that celebrate ecological uniqueness, recognise nature as an active agent in fostering respect, and encourage participatory care for natural spaces. They also highlighted integrating non-human ecosystems, such as pollinator habitats, into human environments to shape behaviour, foster biodiversity appreciation, and challenge cultural assumptions that separate humans from nature. 


\paragraph{TO2 - Limitations and Criticisms:} Experts criticised reliance on visual perception as the primary sensory dimension for design. While impactful and accessible, this overemphasis neglects opportunities to create more immersive multi-sensory experiences. Additionally, experts  pointed out the dilution of biophilic philosophy over time through superficial approaches such as the use of fake plants that prioritise aesthetics over authenticity and fail to reflect natural processes. A lack of ecological coherence was also discussed, with many designs featuring token, isolated elements that fail to provide ecological continuity with the surrounding spaces. 

\paragraph{TO3 - Sensory Design Opportunities:} Experts highlighted the underutilisation of smell, sound, and taste in biophilic design. Smell, despite its emotional and memory-triggering capabilities~\cite{henshaw2013urban}, is often avoided due to concerns about undesirable natural odours. Sound practices focus on noise reduction (also pointed out by \citet{spence2020senses} in his work summarising multisensory design) rather than leveraging natural auditory elements like birdsong to enrich sensory experiences. Taste is rarely addressed, limiting its potential in specific, niche contexts (e.g. ~\cite{horwitz2004eating}). Experts also spoke of opportunities for biophilic spaces that respond to human movement to enhance users' sense of self and foster stronger connections to their environment. Multi-sensory designs that facilitate meditation, immersion, and exploration of nature’s systems (e.g. microbial systems~\cite{risseeuw2024re}) can provoke curiosity.

\paragraph{TO4 - Technology:}
Experts put forth the notion of ``smart nature'' or ``extra nature'', considering technology as a design material for biophilic experiences in architectural spaces. In addition to climate-responsive innovations like hydro-gel windows that regulate sunlight for occupant comfort~\cite{zhou2020liquid}, they proposed technology such as \emph{reactive biophilic booths}, retrofitted with lighting, air quality, and sound sensors, as a modern extension of sun-bathing booths to provide intense, immersive nature-inspired experiences. Experts also discussed leveraging relatively ``low-fi'' technology such as QR codes and mobile applications as digital cues to prompt interaction with natural elements. Generative AI was highlighted for its potential to create innovative visual biophilic designs~\cite{sabinson2023walk, xing2024exploring}, while Mixed Reality (MR) technologies were suggested for their ability to blur artificial and natural environments.

\paragraph{TO5 - Socio-Economic Challenges:} Experts surfaced regulatory, economic, and technical constraints that limit the extent to which biophilic design is implemented in indoor spaces. Reliance on nurseries for specific plant species can pose logistical challenges, while project tender processes often prioritise cost efficiency, stifling opportunities for creative and sustainable solutions. In heritage areas, strict preservation laws and morphological constraints complicate efforts to balance cultural preservation with ecological improvements. Experts also highlighted trade-offs in urban environments where vertical expansion is favoured, compromising ground-level connectivity and ecological continuity.


\paragraph{TO6 - Biophilic Design Values:} Experts discussed the importance of embedding societal values such as neurodiversity~\cite{bruyere2022neurodiversity}, feminism~\cite{boys10there}, and community justice~\cite{fieuw2022towards} into biophilic design for indoor spaces. They also highlighted the need for these environments to cater to marginalised and transient communities that can benefit from nature connectedness for emotional and physical healing. Artistic expression through storytelling and spatial design was identified as a means to reflect community identity (as done by~\cite{samaha2023decolonizing}). By integrating local narratives into natural spaces, biophilic design can strengthen cultural ties and foster deeper personal connections to nature. Experts also noted that collaborative approaches could empower communities to actively care for and sustain these environments.


\paragraph{TO7 - Speculative and Future Visions:}
Experts envisioned future biophilic design that can create integrated ecosystems within indoor spaces, such as biotopes and ecological corridors, to support building wildlife (previously discussed at a city scale~\cite{zielinska2022rethinking}). These participatory ecosystems were seen as fostering species coexistence and offering novel experiences, including moments of discomfort to enrich nature interactions. They also proposed \emph{emotional biophilic spaces} designed to adapt to body rhythms, emotional states, and gender-specific needs. This vision included incorporating natural rhythms, such as tides and seasonal changes, to evoke a sense of belonging to larger natural cycles. Experts also suggested leveraging natural light to create transparency and reflections, evoking concepts like infinite vistas and horizon views.
